% Options for packages loaded elsewhere
\PassOptionsToPackage{unicode}{hyperref}
\PassOptionsToPackage{hyphens}{url}
%
\documentclass[
  11pt,
]{article}
\usepackage{amsmath,amssymb}
\usepackage{iftex}
\ifPDFTeX
  \usepackage[T1]{fontenc}
  \usepackage[utf8]{inputenc}
  \usepackage{textcomp} % provide euro and other symbols
\else % if luatex or xetex
  \usepackage{unicode-math} % this also loads fontspec
  \defaultfontfeatures{Scale=MatchLowercase}
  \defaultfontfeatures[\rmfamily]{Ligatures=TeX,Scale=1}
\fi
\usepackage{lmodern}
\ifPDFTeX\else
  % xetex/luatex font selection
\fi
% Use upquote if available, for straight quotes in verbatim environments
\IfFileExists{upquote.sty}{\usepackage{upquote}}{}
\IfFileExists{microtype.sty}{% use microtype if available
  \usepackage[]{microtype}
  \UseMicrotypeSet[protrusion]{basicmath} % disable protrusion for tt fonts
}{}
\makeatletter
\@ifundefined{KOMAClassName}{% if non-KOMA class
  \IfFileExists{parskip.sty}{%
    \usepackage{parskip}
  }{% else
    \setlength{\parindent}{0pt}
    \setlength{\parskip}{6pt plus 2pt minus 1pt}}
}{% if KOMA class
  \KOMAoptions{parskip=half}}
\makeatother
\usepackage{xcolor}
\usepackage[margin=2.2cm]{geometry}
\usepackage{longtable,booktabs,array}
\usepackage{calc} % for calculating minipage widths
% Correct order of tables after \paragraph or \subparagraph
\usepackage{etoolbox}
\makeatletter
\patchcmd\longtable{\par}{\if@noskipsec\mbox{}\fi\par}{}{}
\makeatother
% Allow footnotes in longtable head/foot
\IfFileExists{footnotehyper.sty}{\usepackage{footnotehyper}}{\usepackage{footnote}}
\makesavenoteenv{longtable}
\setlength{\emergencystretch}{3em} % prevent overfull lines
\providecommand{\tightlist}{%
  \setlength{\itemsep}{0pt}\setlength{\parskip}{0pt}}
\setcounter{secnumdepth}{-\maxdimen} % remove section numbering
\ifLuaTeX
  \usepackage{selnolig}  % disable illegal ligatures
\fi
\IfFileExists{bookmark.sty}{\usepackage{bookmark}}{\usepackage{hyperref}}
\IfFileExists{xurl.sty}{\usepackage{xurl}}{} % add URL line breaks if available
\urlstyle{same}
\hypersetup{
  pdftitle={Worksheet 5A: Angles, Radians, and the Unit Circle},
  hidelinks,
  pdfcreator={LaTeX via pandoc}}

\title{Worksheet 5A: Angles, Radians, and the Unit Circle}
\author{}
\date{}

\begin{document}
\maketitle

\begin{quote}
\textbf{What This Worksheet Is About}

Trigonometry began with ancient astronomers trying to measure the
heavens. Today it's essential for everything from GPS satellites to
semiconductor manufacturing.

We start with two ways to measure angles (degrees and radians) and
introduce the unit circle --- a tool that will make everything else in
trigonometry click into place.
\end{quote}

\begin{center}\rule{0.5\linewidth}{0.5pt}\end{center}

\hypertarget{part-a-degrees-and-radians-two-languages-for-angles-20-problems}{%
\subsection{Part A: Degrees and Radians --- Two Languages for Angles (20
problems)}\label{part-a-degrees-and-radians-two-languages-for-angles-20-problems}}

\hypertarget{why-two-systems}{%
\subsubsection{Why Two Systems?}\label{why-two-systems}}

\textbf{Degrees} come from ancient Babylon (360° in a circle ---
probably because 360 has many divisors and roughly matches days in a
year).

\textbf{Radians} come from mathematics itself: one radian is the angle
where the arc length equals the radius.

\[\boxed{2\pi \text{ radians} = 360°}\]

Therefore: \(\pi \text{ rad} = 180°\), and
\(1 \text{ rad} = \frac{180°}{\pi} \approx 57.3°\)

\begin{center}\rule{0.5\linewidth}{0.5pt}\end{center}

\hypertarget{converting-between-systems}{%
\subsubsection{Converting Between
Systems}\label{converting-between-systems}}

\textbf{Degrees → Radians:} Multiply by \(\frac{\pi}{180}\)

\textbf{Radians → Degrees:} Multiply by \(\frac{180}{\pi}\)

\begin{center}\rule{0.5\linewidth}{0.5pt}\end{center}

\textbf{Convert to radians (leave answers in terms of \(\pi\)):}

\textbf{1.} \(90° ={}\) \_\_\_\_\_\_\_\_\_\_ rad

\textbf{2.} \(180° ={}\) \_\_\_\_\_\_\_\_\_\_ rad

\textbf{3.} \(45° ={}\) \_\_\_\_\_\_\_\_\_\_ rad

\textbf{4.} \(60° ={}\) \_\_\_\_\_\_\_\_\_\_ rad

\textbf{5.} \(30° ={}\) \_\_\_\_\_\_\_\_\_\_ rad

\textbf{6.} \(120° ={}\) \_\_\_\_\_\_\_\_\_\_ rad

\textbf{7.} \(270° ={}\) \_\_\_\_\_\_\_\_\_\_ rad

\textbf{8.} \(315° ={}\) \_\_\_\_\_\_\_\_\_\_ rad

\textbf{9.} \(150° ={}\) \_\_\_\_\_\_\_\_\_\_ rad

\textbf{10.} \(225° ={}\) \_\_\_\_\_\_\_\_\_\_ rad

\begin{center}\rule{0.5\linewidth}{0.5pt}\end{center}

\textbf{Convert to degrees:}

\textbf{11.} \(\frac{\pi}{6}\) rad \(={}\) \_\_\_\_\_\_\_\_\_\_°

\textbf{12.} \(\frac{\pi}{4}\) rad \(={}\) \_\_\_\_\_\_\_\_\_\_°

\textbf{13.} \(\frac{\pi}{3}\) rad \(={}\) \_\_\_\_\_\_\_\_\_\_°

\textbf{14.} \(\frac{2\pi}{3}\) rad \(={}\) \_\_\_\_\_\_\_\_\_\_°

\textbf{15.} \(\frac{5\pi}{6}\) rad \(={}\) \_\_\_\_\_\_\_\_\_\_°

\textbf{16.} \(\frac{5\pi}{4}\) rad \(={}\) \_\_\_\_\_\_\_\_\_\_°

\textbf{17.} \(\frac{7\pi}{6}\) rad \(={}\) \_\_\_\_\_\_\_\_\_\_°

\textbf{18.} \(\frac{11\pi}{6}\) rad \(={}\) \_\_\_\_\_\_\_\_\_\_°

\begin{center}\rule{0.5\linewidth}{0.5pt}\end{center}

** Science Connection: Why Radians Matter**

\textbf{19.} In calculus, the derivative of \(\sin(x)\) is \(\cos(x)\)
--- but \textbf{only if \(x\) is in radians!}

If we used degrees, we'd get:
\(\frac{d}{dx}\sin(x°) = \frac{\pi}{180}\cos(x°)\)

That ugly factor disappears with radians. This is why all scientific
computing uses radians.

\textbf{20.} A satellite dish must be aimed within \(0.1°\) of the
correct angle. Convert this to radians.

\begin{center}\rule{0.5\linewidth}{0.5pt}\end{center}

\hypertarget{part-b-arc-length-and-the-radian-definition-12-problems}{%
\subsection{Part B: Arc Length and the Radian Definition (12
problems)}\label{part-b-arc-length-and-the-radian-definition-12-problems}}

\hypertarget{the-key-formula}{%
\subsubsection{The Key Formula}\label{the-key-formula}}

For a circle of radius \(r\) and central angle \(\theta\) (in radians):

\[\boxed{s = r\theta}\]

where \(s\) is the \textbf{arc length}.

This formula is beautifully simple --- but only works with radians!

\begin{center}\rule{0.5\linewidth}{0.5pt}\end{center}

\textbf{21.} A circle has radius 10 cm. Find the arc length for a
central angle of: - (a) \(\frac{\pi}{2}\) rad - (b) \(\pi\) rad - (c)
\(2\pi\) rad (verify this gives the circumference!)

\textbf{22.} Earth's radius is approximately 6,371 km.

\begin{itemize}
\item
  \begin{enumerate}
  \def\labelenumi{(\alph{enumi})}
  \tightlist
  \item
    If you travel along Earth's surface through an angle of 1 radian (as
    measured from Earth's center), how far have you traveled?
  \end{enumerate}
\item
  \begin{enumerate}
  \def\labelenumi{(\alph{enumi})}
  \setcounter{enumi}{1}
  \tightlist
  \item
    Dublin is at latitude 53.3°N. Convert this to radians.
  \end{enumerate}
\item
  \begin{enumerate}
  \def\labelenumi{(\alph{enumi})}
  \setcounter{enumi}{2}
  \tightlist
  \item
    What is the arc length from the equator to Dublin along a meridian?
  \end{enumerate}
\end{itemize}

\begin{center}\rule{0.5\linewidth}{0.5pt}\end{center}

** Science Connection: How Eratosthenes Measured Earth**

Around 240 BCE, Eratosthenes noticed that at noon on the summer
solstice, the sun was directly overhead in Syene (modern Aswan, Egypt)
--- it shone straight down a well.

But in Alexandria, 800 km north, the sun cast a shadow at an angle of
about 7.2° from vertical.

\textbf{23.} - (a) Convert 7.2° to radians. - (b) The arc length between
the cities is 800 km. Using \(s = r\theta\), solve for \(r\) (Earth's
radius). - (c) The actual radius is about 6,371 km. How close was this
ancient measurement?

\begin{center}\rule{0.5\linewidth}{0.5pt}\end{center}

** Astronomy: Angular Size**

When we look at distant objects, we measure their \textbf{angular
diameter} --- how big they appear in the sky.

\textbf{24.} The Moon's angular diameter is about 0.52° as seen from
Earth. - (a) Convert to radians. - (b) The Moon is about 384,400 km
away. Using \(s = r\theta\) (treating the Moon's diameter as an arc),
estimate the Moon's actual diameter. - (c) The actual diameter is about
3,474 km. How close is your estimate?

\textbf{25.} The Sun's angular diameter is also about 0.53° --- almost
the same as the Moon! This is why solar eclipses are possible. The Sun
is about 150 million km away. Estimate the Sun's diameter.

\textbf{26.} Jupiter's angular diameter varies from 0.55 arcminutes to
0.83 arcminutes depending on its distance from Earth (1 arcminute = 1/60
of a degree). - (a) Convert 0.83 arcminutes to radians. - (b) At closest
approach, Jupiter is about 588 million km away. Estimate its diameter.

\begin{center}\rule{0.5\linewidth}{0.5pt}\end{center}

\hypertarget{part-c-the-unit-circle-your-new-best-friend-24-problems}{%
\subsection{Part C: The Unit Circle --- Your New Best Friend (24
problems)}\label{part-c-the-unit-circle-your-new-best-friend-24-problems}}

\hypertarget{what-is-it}{%
\subsubsection{What Is It?}\label{what-is-it}}

The \textbf{unit circle} is a circle with radius 1, centered at the
origin.

For any angle \(\theta\) measured from the positive x-axis: - The
x-coordinate of the point on the circle is \(\cos(\theta)\) - The
y-coordinate is \(\sin(\theta)\)

\[\boxed{(\cos\theta, \sin\theta) = \text{coordinates on unit circle at angle } \theta}\]

\begin{center}\rule{0.5\linewidth}{0.5pt}\end{center}

\hypertarget{special-angles-to-memorize}{%
\subsubsection{Special Angles to
Memorize}\label{special-angles-to-memorize}}

\begin{longtable}[]{@{}llll@{}}
\toprule\noalign{}
Angle (°) & Angle (rad) & \(\cos\theta\) & \(\sin\theta\) \\
\midrule\noalign{}
\endhead
\bottomrule\noalign{}
\endlastfoot
0° & 0 & 1 & 0 \\
30° & \(\frac{\pi}{6}\) & \(\frac{\sqrt{3}}{2}\) & \(\frac{1}{2}\) \\
45° & \(\frac{\pi}{4}\) & \(\frac{\sqrt{2}}{2}\) &
\(\frac{\sqrt{2}}{2}\) \\
60° & \(\frac{\pi}{3}\) & \(\frac{1}{2}\) & \(\frac{\sqrt{3}}{2}\) \\
90° & \(\frac{\pi}{2}\) & 0 & 1 \\
\end{longtable}

\textbf{Memory trick:} For sine at 0°, 30°, 45°, 60°, 90°:
\(\frac{\sqrt{0}}{2}, \frac{\sqrt{1}}{2}, \frac{\sqrt{2}}{2}, \frac{\sqrt{3}}{2}, \frac{\sqrt{4}}{2}\)

\begin{center}\rule{0.5\linewidth}{0.5pt}\end{center}

\textbf{Find the exact values without a calculator:}

\textbf{27.} \(\cos(0) ={}\) \_\_\_\_\_ , \(\sin(0) ={}\) \_\_\_\_\_

\textbf{28.} \(\cos\left(\frac{\pi}{6}\right) ={}\) \_\_\_\_\_ ,
\(\sin\left(\frac{\pi}{6}\right) ={}\) \_\_\_\_\_

\textbf{29.} \(\cos\left(\frac{\pi}{4}\right) ={}\) \_\_\_\_\_ ,
\(\sin\left(\frac{\pi}{4}\right) ={}\) \_\_\_\_\_

\textbf{30.} \(\cos\left(\frac{\pi}{3}\right) ={}\) \_\_\_\_\_ ,
\(\sin\left(\frac{\pi}{3}\right) ={}\) \_\_\_\_\_

\textbf{31.} \(\cos\left(\frac{\pi}{2}\right) ={}\) \_\_\_\_\_ ,
\(\sin\left(\frac{\pi}{2}\right) ={}\) \_\_\_\_\_

\begin{center}\rule{0.5\linewidth}{0.5pt}\end{center}

\hypertarget{extending-to-all-four-quadrants}{%
\subsubsection{Extending to All Four
Quadrants}\label{extending-to-all-four-quadrants}}

\textbf{ASTC Rule:} Which functions are positive in each quadrant? -
\textbf{A}ll (Quadrant I): Both sin and cos positive - \textbf{S}in
(Quadrant II): Only sin positive - \textbf{T}an (Quadrant III): Only tan
positive - \textbf{C}os (Quadrant IV): Only cos positive

\textbf{Reference Angle:} The acute angle to the x-axis. Use symmetry!

\begin{center}\rule{0.5\linewidth}{0.5pt}\end{center}

\textbf{Find the exact values:}

\textbf{32.} \(\cos\left(\frac{2\pi}{3}\right) ={}\) \_\_\_\_\_ ,
\(\sin\left(\frac{2\pi}{3}\right) ={}\) \_\_\_\_\_

\emph{Hint: \(\frac{2\pi}{3}\) is in Q2, reference angle is
\(\frac{\pi}{3}\)}

\textbf{33.} \(\cos\left(\frac{3\pi}{4}\right) ={}\) \_\_\_\_\_ ,
\(\sin\left(\frac{3\pi}{4}\right) ={}\) \_\_\_\_\_

\textbf{34.} \(\cos\left(\frac{5\pi}{6}\right) ={}\) \_\_\_\_\_ ,
\(\sin\left(\frac{5\pi}{6}\right) ={}\) \_\_\_\_\_

\textbf{35.} \(\cos(\pi) ={}\) \_\_\_\_\_ , \(\sin(\pi) ={}\) \_\_\_\_\_

\textbf{36.} \(\cos\left(\frac{7\pi}{6}\right) ={}\) \_\_\_\_\_ ,
\(\sin\left(\frac{7\pi}{6}\right) ={}\) \_\_\_\_\_

\textbf{37.} \(\cos\left(\frac{5\pi}{4}\right) ={}\) \_\_\_\_\_ ,
\(\sin\left(\frac{5\pi}{4}\right) ={}\) \_\_\_\_\_

\textbf{38.} \(\cos\left(\frac{4\pi}{3}\right) ={}\) \_\_\_\_\_ ,
\(\sin\left(\frac{4\pi}{3}\right) ={}\) \_\_\_\_\_

\textbf{39.} \(\cos\left(\frac{3\pi}{2}\right) ={}\) \_\_\_\_\_ ,
\(\sin\left(\frac{3\pi}{2}\right) ={}\) \_\_\_\_\_

\textbf{40.} \(\cos\left(\frac{5\pi}{3}\right) ={}\) \_\_\_\_\_ ,
\(\sin\left(\frac{5\pi}{3}\right) ={}\) \_\_\_\_\_

\textbf{41.} \(\cos\left(\frac{7\pi}{4}\right) ={}\) \_\_\_\_\_ ,
\(\sin\left(\frac{7\pi}{4}\right) ={}\) \_\_\_\_\_

\textbf{42.} \(\cos\left(\frac{11\pi}{6}\right) ={}\) \_\_\_\_\_ ,
\(\sin\left(\frac{11\pi}{6}\right) ={}\) \_\_\_\_\_

\begin{center}\rule{0.5\linewidth}{0.5pt}\end{center}

\hypertarget{negative-angles-and-angles-beyond-2pi}{%
\subsubsection{\texorpdfstring{Negative Angles and Angles Beyond
\(2\pi\)}{Negative Angles and Angles Beyond 2\textbackslash pi}}\label{negative-angles-and-angles-beyond-2pi}}

\textbf{43.} \(\cos(-\frac{\pi}{4}) ={}\) \_\_\_\_\_ ,
\(\sin(-\frac{\pi}{4}) ={}\) \_\_\_\_\_

\textbf{44.} \(\cos(-\frac{\pi}{6}) ={}\) \_\_\_\_\_ ,
\(\sin(-\frac{\pi}{6}) ={}\) \_\_\_\_\_

\textbf{45.} \(\cos\left(\frac{13\pi}{6}\right) ={}\) \_\_\_\_\_ ,
\(\sin\left(\frac{13\pi}{6}\right) ={}\) \_\_\_\_\_

\emph{Hint: \(\frac{13\pi}{6} = 2\pi + \frac{\pi}{6}\)}

\textbf{46.} \(\cos(3\pi) ={}\) \_\_\_\_\_ , \(\sin(3\pi) ={}\)
\_\_\_\_\_

\begin{center}\rule{0.5\linewidth}{0.5pt}\end{center}

\hypertarget{the-pythagorean-identity}{%
\subsubsection{The Pythagorean
Identity}\label{the-pythagorean-identity}}

\textbf{47.} On the unit circle, a point is at
\((\cos\theta, \sin\theta)\). The distance from this point to the origin
is the radius, which equals 1.

Using the distance formula: \(\sqrt{\cos^2\theta + \sin^2\theta} = 1\)

Square both sides: \(\cos^2\theta + \sin^2\theta ={}\) \_\_\_\_\_

\textbf{This is the most important identity in trigonometry!}

\textbf{48.} Verify the identity for \(\theta = \frac{\pi}{4}\):

\(\cos^2\left(\frac{\pi}{4}\right) + \sin^2\left(\frac{\pi}{4}\right) = \left(\frac{\sqrt{2}}{2}\right)^2 + \left(\frac{\sqrt{2}}{2}\right)^2 ={}\)
\_\_\_\_\_

\textbf{49.} Verify for \(\theta = \frac{\pi}{3}\).

\textbf{50.} If \(\sin\theta = \frac{3}{5}\) and \(\theta\) is in
Quadrant I, find \(\cos\theta\).

\begin{center}\rule{0.5\linewidth}{0.5pt}\end{center}

\hypertarget{part-d-the-tangent-function-and-more-14-problems}{%
\subsection{Part D: The Tangent Function and More (14
problems)}\label{part-d-the-tangent-function-and-more-14-problems}}

\hypertarget{definition}{%
\subsubsection{Definition}\label{definition}}

\[\tan\theta = \frac{\sin\theta}{\cos\theta} = \frac{y}{x} = \text{slope of the line from origin to point on circle}\]

\begin{center}\rule{0.5\linewidth}{0.5pt}\end{center}

\textbf{Find the exact values:}

\textbf{51.} \(\tan(0) ={}\) \_\_\_\_\_

\textbf{52.} \(\tan\left(\frac{\pi}{6}\right) ={}\) \_\_\_\_\_

\textbf{53.} \(\tan\left(\frac{\pi}{4}\right) ={}\) \_\_\_\_\_

\textbf{54.} \(\tan\left(\frac{\pi}{3}\right) ={}\) \_\_\_\_\_

\textbf{55.} \(\tan\left(\frac{\pi}{2}\right) ={}\) \_\_\_\_\_ (What
happens here?)

\textbf{56.} \(\tan\left(\frac{2\pi}{3}\right) ={}\) \_\_\_\_\_

\textbf{57.} \(\tan\left(\frac{3\pi}{4}\right) ={}\) \_\_\_\_\_

\textbf{58.} \(\tan(\pi) ={}\) \_\_\_\_\_

\begin{center}\rule{0.5\linewidth}{0.5pt}\end{center}

** Science Connection: Angles of Incidence and Reflection**

\textbf{59.} In optics, when light hits a surface, the angle of
incidence equals the angle of reflection (measured from the
perpendicular ``normal'' line).

A laser beam hits a mirror at an angle of \(\frac{\pi}{6}\) radians from
the normal. - (a) What is the angle of reflection? - (b) What is the
total deflection of the beam (angle between incoming and outgoing rays)?

\textbf{60.} In a periscope, light bounces off two mirrors. If each
mirror is angled at 45° to the horizontal, through what total angle is
the light path bent?

\begin{center}\rule{0.5\linewidth}{0.5pt}\end{center}

** Microscopy: Numerical Aperture**

The \textbf{numerical aperture} (NA) of a microscope objective
determines its resolving power:

\[NA = n \cdot \sin(\theta)\]

where \(n\) is the refractive index of the medium and \(\theta\) is the
half-angle of the cone of light entering the objective.

\textbf{61.} An air-immersion objective (\(n = 1.0\)) has a half-angle
of 60°. - (a) Find \(\sin(60°)\). - (b) Calculate the numerical
aperture.

\textbf{62.} An oil-immersion objective (\(n = 1.52\)) has a half-angle
of 70°. - (a) Find \(\sin(70°)\) using a calculator (≈ 0.940). - (b)
Calculate the numerical aperture. - (c) Why does using oil allow for
higher NA and better resolution?

\begin{center}\rule{0.5\linewidth}{0.5pt}\end{center}

** Lithography: The Diffraction Limit**

In semiconductor manufacturing, the smallest feature size that can be
printed is limited by:

\[\text{Resolution} = k_1 \cdot \frac{\lambda}{NA}\]

where \(\lambda\) is the wavelength of light, NA is the numerical
aperture, and \(k_1\) is a process factor (typically around 0.25-0.4).

\textbf{63.} Extreme ultraviolet (EUV) lithography uses
\(\lambda = 13.5\) nm and has NA ≈ 0.33. - (a) With \(k_1 = 0.3\), what
is the theoretical minimum feature size? - (b) Current chips have 3nm
features. Does this seem achievable with these parameters?

\textbf{64.} Older lithography used deep ultraviolet (DUV) light at
\(\lambda = 193\) nm with NA ≈ 1.35 (using immersion). - (a) With
\(k_1 = 0.3\), what resolution is achievable? - (b) Why did the industry
need to move to EUV for smaller features?

\begin{center}\rule{0.5\linewidth}{0.5pt}\end{center}

\end{document}
